\addcontentsline{toc}{subsubsection}{Solution 33. The Dyadic Parallelogram--Quadratic Squeeze}

\begin{dlsummary}[title={\bfseries Solution 33. The Dyadic Parallelogram--Quadratic Squeeze},breakable]

\begin{proof}
\textbf{Step 1. Obtain an inequality for the residual.}

Put $g(x)=f(x)-x^2$. Comparing the outer expressions in the given chain
and substituting $f(x)=x^2+g(x)$ gives
\[
2g(x)+3g(y)+2(x-2y)^2\geq0
\qquad(x,y\in\mathbb R). \tag{1}
\]
Taking $x=y=0$ shows $g(0)\geq0$, while setting $x=2t$ and $y=t$ gives
\[
2g(2t)+3g(t)\geq0\qquad(t\in\mathbb R). \tag{2}
\]

\medskip
\noindent
\textbf{Step 2. Prove that $g$ is nonnegative.}

The right-hand inequality in the problem and the parallelogram identity
imply
\[
g(x+y)+g(x-y)\leq2g(x)
\qquad(x,y\in\mathbb R). \tag{3}
\]
Taking $x=y=t$ gives $g(2t)\leq2g(t)-g(0)$. Combining this with (2),
\[
0\leq2g(2t)+3g(t)\leq7g(t)-2g(0).
\]
Since $g(0)\geq0$, it follows that $g(t)\geq\frac27g(0)\geq0$ for
all $t\in\mathbb R$.

\medskip
\noindent
\textbf{Step 3. Show that $g$ is constant.}

Fix $a,b\in\mathbb R$, and define
\[
x_k=a+k(b-a)\quad(k\geq0),
\qquad
\Delta_k=g(x_{k+1})-g(x_k).
\]
For each $k\geq1$, apply (3) with $x=x_k$ and $y=b-a$. Then
\[
g(x_{k+1})+g(x_{k-1})\leq2g(x_k),
\]
so $\Delta_k\leq\Delta_{k-1}$. Therefore
$\Delta_k\leq\Delta_0=g(b)-g(a)$ for all $k\geq0$. For any positive
integer $n$, nonnegativity of $g$ yields
\[
-g(a)\leq g(x_n)-g(a)
=\sum_{k=0}^{n-1}\Delta_k
\leq n\bigl(g(b)-g(a)\bigr).
\]
Letting $n\to\infty$ gives $g(b)\geq g(a)$. Interchanging $a$ and $b$
shows that $g$ is constant, say $g(x)=d$ with $d\geq0$. Hence
\[
f(x)=x^2+d. \tag{4}
\]

\medskip
\noindent
\textbf{Step 4. Determine $d$ and verify the solution.}

Using (4) in the composition condition gives
\[
1\geq f(f(1))=(1+d)^2+d=1+3d+d^2.
\]
As $d\geq0$, this forces $d=0$, so $f(x)=x^2$ for all real $x$.

For $f(x)=x^2$, the middle and right-hand expressions in the problem
both equal $5x^2+5y^2$, while the left-hand expression is
\[
3x^2+8xy-3y^2=5x^2+5y^2-2(x-2y)^2
\leq5x^2+5y^2.
\]
Also $f(f(1))=1$. Thus $f(x)=x^2$ satisfies all conditions and is the
unique solution.
\end{proof}

\end{dlsummary}